% !TeX root = ../se200_identity_regimes.tex

\section{Related Work}
\label{sec:related}

This paper is situated at the intersection of formal ontology, knowledge
representation, and the study of classification under disagreement.
Rather than surveying ontology engineering broadly, this section situates the
present result relative to prior work on
classification under disagreement, identity and roles in formal ontology,
and upper-ontology design.

\subsection{Classification, Disagreement, and Neutrality}

A substantial body of work has shown that classification systems embed
institutional, social, and political commitments rather than functioning
as neutral descriptive devices~\citep{bowker1999sorting,haslanger2012resisting}.
In scientific and policy contexts, disagreement is often persistent rather
than transient, reflecting incompatible evidentiary, methodological, or
evaluative frameworks rather than resolvable empirical uncertainty~\citep{longino1990science}.

These observations motivate the need for representations that remain usable
without assuming interpretive convergence.
The present work adopts this motivation but addresses a distinct structural
question: which referential-regime distinctions become unavoidable when
disagreement is treated as persistent and causal or normative commitments
are externalized from the substrate rather than settled within it.
The answer given here is conditional on the neutrality, stable-reference,
and no-hidden-regime assumptions stated above: under those assumptions,
the analysis yields a current lower bound on the referential-regime structure
required for accountability under persistent disagreement.


\subsection{Identity, Roles, and Ontological Methodology}

Formal ontology has long emphasized the role of identity conditions and
persistence criteria in category formation.
The OntoClean methodology~\citep{guarino2002evaluating} clarifies why roles are
anti-rigid and unsuitable as identity-defining categories~\citep{guarino1998roles},
and subsequent work on social and institutional entities has further
developed these distinctions~\citep{masolo2003wonderweb}.

This paper builds on these insights but applies them at the substrate level.
The present analysis operates at a different level than role-based modeling,
deriving substrate-level referential regimes
from neutrality requirements rather than from domain classifications.

\subsection{Upper Ontologies}

Upper ontologies such as BFO~\citep{smith2015basic}, DOLCE~\citep{gangemi2002dolce},
and UFO~\citep{guizzardi2005ontological} provide rich accounts of dependence,
endurants and perdurants, and social objects.
They differ substantially in commitments and optimization objectives, and
are not designed primarily as neutral substrates for interoperability under
persistent disagreement.

The present work does not compare expressivity or advocate alignment with any
particular upper ontology.
Instead, it derives lower-bound constraints that apply independently of such
frameworks, conditional on the neutrality and stability requirements adopted here.
The regime structure derived here can be situated with respect to two
widely-used frameworks to clarify where under-constraints arise
not as defects in those frameworks but as points where implementations
built on them would need to introduce additional commitments
to satisfy the neutrality and stability requirements assumed here.

BFO~\citep{smith2015basic} organizes entities under the
continuant/occurrent distinction:
continuants persist through time
while occurrents unfold across it~\citep{arp2015building}.
OCC (time-indexed occurrences) maps cleanly onto BFO occurrents.
ENR through NOR map onto independent and generically dependent continuants,
but BFO does not enforce the regime-level separation between them:
obligation-bearing entities (OBL), authority-bearing structures (NOR),
and descriptive indicators (REC) are all candidates for
Generically Dependent Continuant, leaving their
stability-critical distinctions to the implementer.
CTX (applicability and scope) has no first-class BFO home;
non-spatial applicability contexts such as normative jurisdiction
or conditional scope are represented as GDC, effectively
conflating CTX with NOR under a BFO-based implementation.
Under persistent disagreement, these conflations require
hidden discriminators to recover regime-level distinctions,
introducing implicit structure of the kind
prohibited by the neutrality constraint.

DOLCE~\citep{gangemi2002dolce} introduces social objects as
first-class entities and provides a richer treatment of
institutional and normative structure than BFO~\citep{masolo2003wonderweb}.
This makes DOLCE a better fit for OBL and NOR as distinct categories.
However, the NOR/REC conflation persists: both normative structures and
descriptive indicators are candidates for social object or quality,
and the distinction between authority-as-referent and
descriptive-record-as-referent is not enforced at the type level.
DOLCE's region-based account of scope is spatiotemporal,
leaving normative applicability contexts (CTX) without a
structural home distinct from NOR or REC,
and the OBL/NOR boundary (between obligation-bearer and
normative structure) is mediated by participation relations
rather than enforced by regime membership~\citep{masolo2003wonderweb}.
As with BFO, implementations resolving these under-constraints
under persistent disagreement silently introduce hidden regimes.

These observations do not constitute objections to BFO or DOLCE
as general-purpose ontological frameworks.
They locate a specific point where the optimization objectives
of those frameworks (expressiveness, metaphysical coverage,
alignment with scientific practice) diverge from the
optimization objective assumed here,
and where that divergence has structural consequences
for neutral interoperability under persistent disagreement.

When constrained to enforce conservative extension
over identity and admissibility conditions,
such frameworks can realize the regime structure derived here.
The contribution of this paper is to identify the structural constraints
required when neutrality under persistent disagreement is taken as the primary objective.

The present result is therefore not a competing upper ontology,
but a constraint on any ontology intended to function as a
neutral, stability-preserving substrate under persistent disagreement.

\subsection{Relation to Common Logic and Ontology Exchange}

The neutral substrate and derived regimes can be
expressed in ISO/IEC 24707 Common Logic~\citep{iso24707} as a conservative theory,
supporting interpreter-independent ontology exchange
without reliance on description-logic restrictions or fixed signatures.
Here, conservative means that
extensions add structure
without altering the consequences
derivable about substrate-level terms~\citep{hodges1997shorter},
ensuring that incompatible interpretive frameworks
do not retroactively change the meaning of shared references.
This enables compatibility with ontology exchange workflows
widely used in applied ontology~\citep{kutz2010carnap}
and digital twin infrastructures~\citep{grieves2016digital},
where multiple interpretive layers must interoperate over stable referents.

\subsection{Causality, Explanation, and Infrastructure}

Separating descriptive structure from causal or explanatory claims is a
well-established principle in causal analysis~\citep{pearl2009causality}.
Similarly, research on large-scale information infrastructures emphasizes
the importance of long-term stability and interpretive restraint across
institutional boundaries~\citep{edwards2011infrastructure}.

These perspectives align with the design neutrality constraints assumed in this paper.
The ontology analyzed here is not an explanatory theory, but an
infrastructural substrate intended to support explanation, evaluation, and
policy analysis without embedding them.

\subsubsection{Relation to FAIR Data Principles}

The FAIR principles promote data reuse across communities and over time,
but do not specify the ontological constraints required for interoperability
under persistent disagreement.
In practice, reuse is often pursued through
negotiated vocabularies or shared interpretations,
which presuppose forms of convergence that may not be
available in legal, political, or analytic settings.

The results of this paper clarify a structural precondition for FAIR-aligned
reuse in pluralistic or institutionally heterogeneous contexts:
stable reference must be supported independently of
causal or normative agreement.
The referential regimes derived here therefore constrain the
ontological substrates upon which FAIR-compliant systems implicitly rely
when reuse occurs under persistent disagreement.


%=====================================

This work is situated at the intersection of ontology,
knowledge representation, and data modeling,
with a focus on identity, persistence, and reference under disagreement.

\subsection{Identity and Persistence}

Philosophical and formal treatments of identity have long distinguished
criteria for persistence across time and change.
Approaches based on endurantism and perdurantism
provide different accounts of how entities persist,
but typically assume a fixed interpretive framework
within which identity conditions are evaluated.

In contrast, this work considers identity conditions that must remain
stable under persistent interpretive disagreement.
The requirement that identity be invariant under admissible transformations
introduces constraints that are not captured by standard accounts
of persistence.

\subsection{Ontology and Knowledge Representation}

In ontology engineering and knowledge representation,
identity is often tied to classification, typing, or role assignment,
as in description logics and related frameworks.
Such systems typically allow identity conditions to depend on
context, interpretation, or modeling choices.

The present work departs from these approaches by requiring that
identity conditions be independent of causal and normative commitments
at the substrate level.
Rather than constructing representations,
it derives structural constraints that any substrate must satisfy
to support neutral reference across interpretive frameworks.

\subsection{Data Modeling and Provenance}

Data models and provenance frameworks track entities across transformations,
including versioning, aggregation, and derivation.
Standards such as the W3C PROV model explicitly represent derivation
and lineage relations, while temporal and versioned data models
capture change over time.

These systems provide mechanisms for tracking entities under transformation,
but may allow identity to depend on representation,
derivation history, or application-specific semantics.

This work abstracts from such domain-specific considerations
and instead identifies transformation-invariant identity regimes
required for stable reference.
The transformation families considered here provide a uniform basis
for analyzing identity behavior across domains.

\subsection{Position of This Work}

Unlike prior work that assumes a fixed interpretive setting,
this paper derives referential regimes from
invariance under admissible transformations in a neutral substrate.
The result is a necessary lower bound on the number of regimes required
to maintain stable reference under persistent interpretive disagreement.
No claim is made of completeness or minimality.
